%% ===================================================================
%% KAPITOLA 5: TESTOVÁNÍ A VYHODNOCENÍ
%% ===================================================================
\section{Testování a vyhodnocení}
\label{sec:testovani}

\subsection{Testovací metodika}

Testování kompilátoru bylo rozděleno do dvou úrovní:
\begin{itemize}
	\item \textbf{jednotkové testy} ověřují jednotlivé moduly (lexer, parser,
	      sémantická analýza, IR, optimalizace a~codegen) na malých izolovaných
	      případech,
	\item \textbf{integrační testy} ověřují celý překladový řetězec přes
	      veřejné CLI rozhraní a~kontrolují, že je vygenerovaný kód vytvořen
	      a~obsahuje očekávané konstrukce.
\end{itemize}

Jednotkové testy jsou implementovány ve dvou formách:
(1) testy přímo v modulech v~\texttt{src/} pomocí \texttt{\#[cfg(test)]} a
(2) testy v~adresáři \texttt{tests/}, kde jsou jednotlivé části kompilátoru
testovány přes veřejné API knihovny \texttt{roth}.

Integrační testy (\texttt{tests/integration\_tests.rs}) spouštějí binární
program přes \texttt{cargo run} a~ověřují vznik výstupu v~adresáři \texttt{.build/}.
Tím se simuluje běžné použití kompilátoru z~příkazové řádky.

Všechny testy lze spustit příkazem:
\begin{lstlisting}[language=bash, caption={Spuštění testů}]
cargo test
\end{lstlisting}

\subsection{Výsledky testování}

Testovací sada pokrývá základní syntaktické konstrukce jazyka, práci s~datovým
zásobníkem, překlad uživatelských definic, vybrané řídicí struktury a~optimalizační
průchody nad IR.

V~adresáři \texttt{tests/} jsou testy rozděleny do tematických souborů:
\texttt{lexer\_tests.rs}, \texttt{parser\_tests.rs}, \texttt{analyzer\_tests.rs},
\texttt{ir\_tests.rs}, \texttt{codegen\_tests.rs} a~\texttt{integration\_tests.rs}.
Tyto testy společně ověřují správnost výstupů jednotlivých fází i~to, že
generování kódu odpovídá očekávané sémantice.

Při vývoji byly odhaleny a~opraveny zejména chyby související s~hranami
tokenizace (pozice tokenů přes víceřádkový vstup, escapování řetězců) a~se
správným zacházením s~kontextem REPL (registrace slov a~proměnných mezi vstupy).

\subsection{Porovnání s existujícími implementacemi}

Cílem práce není nahrazovat zavedené implementace jako Gforth, ale popsat
a~ověřit architekturu kompilátoru založenou na explicitním mezikódu,
optimalizačních průchodech a~více backendech.

Ve srovnání s~tradičním \Forth{} je \\Roth{} koncipován především jako
kompilátorový projekt: jazyková sada je omezena na konstrukce potřebné pro
demonstraci frontendu, middle-endu i~backendů. Těžiště práce proto neleží
v~úplné kompatibilitě se standardem \Forth{}, ale v~návrhu překladového řetězce.
